\begin{table*}[t]
\centering
\caption{E4: guardrail efficacy over 137 proposed agent actions (84 faulty, 53 valid).}
\label{tab:e4-tiers}
\footnotesize
\resizebox{\ifdim\width>\linewidth\linewidth\else\width\fi}{!}{%
\begin{tabular}{llccccrc}
\toprule
Tier & Semantic assets available & Caught & Recall & 95\% CI & Correct diag. & False alarms & FP rate \\
\midrule
G0 & none & 0/84 & 0.000 & [0.00, 0.04] & 0.000 & 0 & 0.000 \\
G1 & JSON call schema (keys, datatypes, enumerations) & 12/84 & 0.143 & [0.08, 0.23] & 0.143 & 0 & 0.000 \\
G2 & tag dictionary, engineering ranges, unit \textbackslash{}emph{string} & 57/84 & 0.679 & [0.57, 0.77] & 0.631 & 4 & 0.075 \\
G3 & G2 + RDF/OWL graph, SHACL, QUDT dimensions and conversions, ISA-95 mereology & 62/84 & 0.738 & [0.64, 0.82] & 0.738 & 0 & 0.000 \\
G4 & G3 + PackML behaviour, explicit interlocks, SOSA/PROV observations & 84/84 & 1.000 & [0.96, 1.00] & 1.000 & 0 & 0.000 \\
\bottomrule
\end{tabular}}
\\[4pt]\begin{minipage}{\linewidth}\footnotesize \emph{Correct diagnosis} requires the validator to name the right fault family, not merely to reject. G2 rejects unit-scale errors for the wrong reason and rejects legal re-expressions of a value in a different unit.\end{minipage}
\end{table*}
